\documentclass[a4paper,10pt,fleqn]{article}
\usepackage[margin=1in]{geometry}
\usepackage{amssymb}
\usepackage{adjustbox}
\usepackage{natbib}
\usepackage[colorlinks=true,linkcolor=blue,citecolor=blue,urlcolor=blue]{hyperref}
\usepackage{fuzz}
\usepackage{zed-maths}
\usepackage{zed-proof}
\newdimen\savedleftskip
\begin{document}

\section*{Simple Schema Reference}

\noindent A horizontal definition gives a new name to an existing schema or schema expression without creating a new boxed schema paragraph.

\bigskip

\begin{schema}{Counter}
n : \nat
\end{schema}

\noindent Plain alias: CounterRef is definitionally equal to Counter.

\bigskip

\begin{zed}
CounterRef \defs Counter
\end{zed}

\noindent Inline schema text: a compact one-line schema form.

\bigskip

\begin{zed}
NatPair \defs [ x : \nat; y : \nat | x < y ]
\end{zed}

\section*{Multi-predicate Inline Text}

\noindent Multiple predicates within the inline text are separated by semicolons and joined with logical-and in the output.

\bigskip

\begin{zed}
BoundedNat \defs [ n : \nat | n > 0 \land n < 100 ]
\end{zed}

\section*{Generic LHS}

\noindent The LHS may carry generic type parameters for parameterised horizontal definitions per Z RM §3.8.

\bigskip

\begin{schema}{GenStack}[X]
top : X
\end{schema}

\begin{zed}
StackAlias[X] \defs GenStack[X]
\end{zed}

\end{document}